\documentclass[12pt,a4paper]{extarticle}
\usepackage{fontspec}
\usepackage{amsmath}
\usepackage{geometry}
\geometry{left=25mm,right=15mm,top=20mm,bottom=20mm}
\usepackage{setspace}
\usepackage{array}
\usepackage{longtable}
\usepackage{booktabs}
\usepackage{ragged2e}
\usepackage{enumitem}
\usepackage{xcolor}

\setmainfont{Times New Roman}
\onehalfspacing
\setlength{\parindent}{1.25em}
\sloppy
\tolerance=2000
\emergencystretch=3em
\pagestyle{plain}
\setcounter{tocdepth}{1}
\renewcommand{\contentsname}{Содержание}

\newcommand{\code}[1]{\texttt{#1}}
\newcommand{\hh}[1]{\small\textbf{#1}}

\begin{document}

%======================= TITLE PAGE =======================
\thispagestyle{empty}
\begin{center}
\vspace*{1.8cm}
{\large Проект \code{physics\_md\_al\_fe}}\\[0.3cm]
{\large Научно-технический отчёт о молекулярно-динамическом моделировании}\\[2.0cm]

{\LARGE\bfseries Передача напряжения от магнитострикционного включения Fe$_4$Al$_{13}$ в алюминиевую матрицу: локальная дефектность у границы без развитой пластической зоны}\\[2.0cm]

{\large Атомистическое моделирование системы «алюминий + интерметаллидное включение Fe$_4$Al$_{13}$»}\\[1.4cm]

{\large Метод: классическая молекулярная динамика\\(программный комплекс LAMMPS, потенциал MEAM, ускорение на видеокарте)}\\[2.4cm]

\vfill
{\large Расчёты выполнил студент: Михайловский Илья Евгеньевич, группа 251-931}\\[0.4cm]
{\large Преподаватель: Пшонкин Даниил Евгеньевич}\\[0.8cm]
{\large 25 июня 2026 г.}
\end{center}
\newpage

%======================= TOC =======================
\tableofcontents
\newpage

%======================= 1 =======================
\section{Краткий вывод (на одну страницу)}

\noindent\textbf{Что сделано.} Выполнены и доведены до физически корректного состояния два молекулярно-динамических расчёта системы «однородная алюминиевая матрица + одно включение Fe$_4$Al$_{13}$». Включение испытывает малую собственную деформацию вдоль оси $Z$, которая служит механическим эквивалентом магнитострикционного воздействия (магнитное поле в расчёт напрямую не вводится). Проверялось главное: передаётся ли напряжение от включения в матрицу и возникают ли при этом дефекты, дислокации или пластическая перестройка.

\smallskip
\noindent\textbf{Что найдено.}
\begin{itemize}[leftmargin=1.4em,itemsep=2pt,topsep=2pt]
\item Включение действительно передаёт напряжение в матрицу. Наиболее нагружена узкая приграничная область — оболочка 0--5~Å у границы включение--матрица. Здесь эквивалентное напряжение по Мизесу достигает 368.194~МПа (расчёт на 254055 атомов), что выше характерного предела текучести алюминия около 120~МПа.
\item У границы возникает локальная дефектная перестройка: всего в матрице насчитывается 4265 дефектных атомов, и практически все они сосредоточены в тонком приграничном слое.
\item В расчёте на 510375 атомов дополнительно обнаружен короткий зародыш частичной дислокации: один дислокационный сегмент суммарной длиной 8.47~Å, тип вектора смещения $\tfrac{1}{6}\langle 112\rangle$. В контрольном расчёте без деформации включения дислокаций не было.
\end{itemize}

\smallskip
\noindent\textbf{Что не найдено.}
\begin{itemize}[leftmargin=1.4em,itemsep=2pt,topsep=2pt]
\item Развитая дислокационная сеть не сформировалась ни в одном расчёте.
\item Развитая пластическая зона в объёме матрицы не образовалась: дефекты не уходят вглубь матрицы. Уже на расстоянии более 15~Å от границы дефектность практически исчезает, а дальняя матрица остаётся почти идеальной.
\item В длинном расчёте на 254055 атомов оформленной дислокационной линии не зафиксировано (0 сегментов).
\end{itemize}

\smallskip
\noindent\textbf{Главный физический вывод.} Гипотеза о локальном механическом воздействии магнитострикционного включения на алюминиевую матрицу подтверждается: включение нагружает матрицу и создаёт у границы физические условия для зарождения дефекта (локальное напряжение выше предела текучести, локальная дефектная перестройка, а в одном из расчётов — короткий зародыш частичной дислокации). При этом развитая пластичность пока не доказана: эффект остаётся локальным, приграничным и не превращается в развитую пластическую зону в объёме матрицы на исследованном режиме.

%======================= 2 =======================
\section{Постановка задачи}
Исследуется механизм передачи напряжения от интерметаллидного включения Fe$_4$Al$_{13}$ в окружающую алюминиевую матрицу. Постановка соответствует последней формулировке преподавателя:
\begin{itemize}[leftmargin=1.4em,itemsep=2pt,topsep=2pt]
\item среда — однородная алюминиевая матрица;
\item в матрице находится \textbf{одно} включение Fe$_4$Al$_{13}$;
\item \textbf{границу зерна специально убрали}: в последней постановке требовалось изучать именно границу «включение--матрица», а не границу зерна, поэтому модель сделана однородной (без межзёренной границы и без искусственно введённых вакансий);
\item магнитострикционное воздействие задаётся механически — как малая собственная деформация атомов включения вдоль оси $Z$;
\item магнитное поле в молекулярную динамику напрямую \textbf{не вводится};
\item проверяется, переходит ли напряжение от включения в матрицу и вызывает ли оно дефекты, дислокации или пластическую перестройку, с особым вниманием к зонам максимального напряжения вдоль оси $Z$.
\end{itemize}
Критерий «подтверждённой дислокации», принятый в расчётной части: число дислокационных сегментов больше нуля, либо суммарная длина дислокационной линии больше нуля, либо плотность дислокаций больше нуля.

%======================= 3 =======================
\section{Физическая модель}
\noindent\textbf{Материалы.} Матрица — алюминий с кубической гранецентрированной решёткой (параметр решётки $a = 4.05$~Å). Включение — интерметаллид Fe$_4$Al$_{13}$ (моноклинная фаза, пространственная группа C2/m). Температура моделирования — 300~К.

\smallskip
\noindent\textbf{Магнитное воздействие и его механический эквивалент.} В проектной постановке рассматривается магнитное поле порядка 0.7~Тл. Поле в расчёте не моделируется явно. Вместо него задаётся собственная (внутренняя) деформация включения вдоль оси $Z$. Связь воздействия с деформацией оценивается так:
\begin{itemize}[leftmargin=1.4em,itemsep=2pt,topsep=2pt]
\item оценочное магнитострикционное напряжение во включении: $\sigma \approx 147$~МПа;
\item модуль упругости алюминия: $E \approx 75.7$~ГПа;
\item отсюда относительная деформация $\varepsilon = \sigma / E = 147~\text{МПа} / 75.7~\text{ГПа} = \mathbf{0.001942}$.
\end{itemize}
Таким образом, выбранное значение деформации 0.001942 — это \textbf{не произвольное число}, а физически обоснованное малое относительное удлинение (примерно 0.19~\%), отвечающее оценочному магнитострикционному напряжению 147~МПа в алюминии. Для сравнения: характерный предел текучести алюминия — около 120~МПа.

\smallskip
\noindent\textbf{Что мы ожидаем увидеть.} Если включение «раздувается» вдоль оси $Z$, оно давит на матрицу. Самая нагруженная область — тонкий слой матрицы у границы включения. Вопрос в том, достаточно ли этого локального напряжения, чтобы в матрице зародились дефекты и дислокации, и насколько далеко такая перестройка распространяется.

%======================= 4 =======================
\section{Методика расчёта}
\noindent\textbf{Метод.} Классическая молекулярная динамика: рассчитывается движение каждого атома под действием межатомных сил; по итоговой атомной конфигурации анализируется структура и локальные напряжения.

\smallskip
\noindent\textbf{Программный комплекс и потенциал.} Расчёты выполнены в программном комплексе LAMMPS (сборка с ускорением на видеокарте, пакет KOKKOS/CUDA; использовалась видеокарта NVIDIA GeForce RTX 3060, 12~ГБ). Межатомное взаимодействие описывается потенциалом семейства MEAM (модифицированный метод внедрённого атома, трек \code{MEAM\_Jelinek\_2012}), пригодным для системы алюминий--железо. Шаг интегрирования по времени в производственной стадии — 0.0002~пс.

\smallskip
\noindent\textbf{Как вводилась деформация.} Сначала система плавно прогревалась от 10~К до 300~К и стабилизировалась при 300~К. Затем атомам включения задавалась малая собственная деформация вдоль оси $Z$ с относительным удлинением 0.001942 (объёмосохраняющее воздействие, приложенное только к включению). После этого система рассчитывалась в производственном режиме, и анализировалось, как напряжение от включения распределяется по матрице.

\smallskip
\noindent\textbf{Какие анализы структуры применялись} (английские названия методов и принятые сокращения вынесены в приложение, раздел~14.3):
\begin{enumerate}[leftmargin=1.6em,itemsep=3pt,topsep=2pt]
\item \textbf{Анализ дислокаций по атомной конфигурации.} По итоговому положению атомов автоматически выделяются дислокационные линии, определяются их число, суммарная длина, плотность и тип (вектор смещения, характеризующий дислокацию).
\item \textbf{Анализ локального кристаллического окружения.} Каждый атом классифицируется по форме его ближайшего окружения: нормальная кубическая гранецентрированная структура (типичная для алюминия); окружение, характерное для гексагональной плотноупакованной структуры (признак дефекта упаковки или следа частичной дислокации); нераспознанное или дефектное окружение.
\item \textbf{Сопоставление локальной атомной структуры с шаблонами кристаллических решёток.} Независимая проверка структуры по эталонным шаблонам решёток; используется как контроль к предыдущему методу.
\item \textbf{Оценка локальных напряжений по зонам.} Локальное напряжение оценивается по вириальному тензору (по покомпонентным напряжениям отдельных атомов) с приведением к зонам матрицы на разном расстоянии от границы включения. Приводятся осевая компонента $P_{zz}$ (вдоль оси воздействия) и эквивалентное напряжение по Мизесу (скалярная мера сдвиговой нагрузки, по которой обычно судят о близости к пластическому течению).
\end{enumerate}

%======================= 5 =======================
\section{Проверка корректности расчёта}
Прежде чем делать физические выводы, проверена техническая достоверность расчётов.

\smallskip
\noindent\textbf{Длинный расчёт на 254055 атомов:}
\begin{itemize}[leftmargin=1.4em,itemsep=2pt,topsep=2pt]
\item расчёт завершён штатно, без ошибок (производственная стадия и анализ завершены успешно, коды возврата нулевые);
\item пройдено 120000 шагов из 120000 запланированных (примерно 24~пс модельного времени);
\item финальная температура 288.39766~К, максимальная за расчёт 291.552~К — то есть система всё время находилась около заданных 300~К, \textbf{перегрева нет};
\item численного разрушения модели нет: атомы не теряются, нет аномального роста температуры или давления;
\item размер системы: расчётная ячейка $14.175 \times 14.175 \times 21.06$~нм, всего 254055 атомов.
\end{itemize}

\smallskip
\noindent\textbf{Расчёт на 510375 атомов:} обе стадии (контроль и физический случай) завершены штатно с нулевыми кодами возврата, максимальная температура 291.98355~К (ниже аварийного порога 1000~К), все необходимые выходные файлы получены, результат признан физически валидным.

\smallskip
\noindent Вывод: оба расчёта технически корректны, их результаты можно интерпретировать физически.

%======================= 6 =======================
\section{Результаты расчёта на 510375 атомов}
Этот расчёт (валидный, на большой системе) включал два случая на одной геометрии: \textbf{контрольный} — без деформации включения (деформация 0); \textbf{физический} — с деформацией включения 0.001942.

\smallskip
Основные результаты:
\begin{itemize}[leftmargin=1.4em,itemsep=2pt,topsep=2pt]
\item число атомов: 510375; деформация в физическом случае: 0.001942; максимальная температура: 291.98355~К;
\item \textbf{контрольный случай}: дислокаций нет (0 сегментов);
\item \textbf{физический случай}: найден \textbf{1 дислокационный сегмент}, суммарная длина дислокационной линии \textbf{8.47~Å}, тип вектора смещения $\tfrac{1}{6}\langle 112\rangle$ (частичная дислокация Шокли).
\end{itemize}
Это короткий \textbf{зародыш частичной дислокации}, а не развитая дислокационная сеть. У границы включение--матрица наблюдались локальные дефекты и повышенные напряжения, однако развитой пластической зоны вглубь матрицы не возникло: по подсчёту дефектов и по низкому эквивалентному напряжению по Мизесу дальняя матрица остаётся практически упругой. Существенно, что сегмент появился именно в физическом случае и отсутствовал в контрольном, — это связывает зародыш дислокации именно с воздействием включения.

%======================= 7 =======================
\section{Результаты длинного расчёта на 254055 атомов (120000 шагов)}
Это основной длинный расчёт, доведённый до 120000 шагов на той же физической постановке.

\smallskip
\noindent\textbf{Параметры расчёта:}
\begin{itemize}[leftmargin=1.4em,itemsep=2pt,topsep=2pt]
\item запуск: \code{runs/stageE\_250k\_single\_physical\_longrun/20260623-205439}, случай \code{E3\_phys001942\_250k\_120k};
\item число атомов: 254055; расчётная ячейка: $14.175 \times 14.175 \times 21.06$~нм;
\item деформация включения: 0.001942; длительность: 120000 шагов;
\item финальная температура 288.39766~К, максимальная 291.552~К.
\end{itemize}

\smallskip
\noindent\textbf{Дислокации.} Оформленных дислокационных линий \textbf{не найдено}: число дислокационных сегментов — \textbf{0}; суммарная длина дислокаций — \textbf{0.0~Å}; плотность дислокаций — \textbf{0.0~м$^{-2}$}.

Это нужно понимать аккуратно. Нулевой результат — это \textbf{не отрицание физического эффекта}, а указание на то, что в выбранном объёме и режиме развитая (оформленная) дислокационная линия не сформировалась. Дефектная перестройка у границы при этом присутствует (см. ниже).

\smallskip
\noindent\textbf{Локальные дефекты.} По анализу локального кристаллического окружения в алюминиевой матрице (всего 243669 атомов матрицы):
\begin{itemize}[leftmargin=1.4em,itemsep=2pt,topsep=2pt]
\item атомов с нормальной кубической гранецентрированной структурой: 239404 — это \textbf{98.25~\%} матрицы;
\item атомов с окружением, характерным для гексагональной плотноупакованной структуры: 7;
\item атомов с нераспознанным или дефектным окружением: 4258;
\item всего дефектных атомов в матрице: \textbf{4265} (4258 + 7);
\item из них за пределами ближней интерфейсной оболочки: всего \textbf{2} дефектных атома;
\item атомов с гексагональным окружением за пределами интерфейсной оболочки: \textbf{0}.
\end{itemize}
Вывод: дефекты локализованы у интерфейса «включение--матрица» и практически \textbf{не распространяются} в объём матрицы.

%======================= 8 =======================
\section{Таблица радиального профиля дефектов и напряжений}
Таблица показывает, как меняются дефектность и напряжение по мере удаления от границы включения (длинный расчёт на 254055 атомов).

\smallskip
\begin{center}
\small
\setlength{\tabcolsep}{6pt}
\renewcommand{\arraystretch}{1.25}
\begin{tabular}{l r r r r r}
\toprule
\hh{Оболочка}     & \hh{Число}  & \shortstack[r]{\hh{Атомы с гекс.}\\\hh{окружением}} & \shortstack[r]{\hh{Атомы с дефектн.}\\\hh{окружением}} & \shortstack[r]{$\mathbf{P_{zz}}$,\\\hh{МПа}} & \shortstack[r]{\hh{По Мизесу,}\\\hh{МПа}} \\
\hh{от границы}    & \hh{атомов} & & & & \\
\midrule
0--5~Å        & 3005   & 6 & 2894 & $-582.896$ & 368.194 \\
5--15~Å       & 13531  & 1 & 1362 & 837.326    & 57.388  \\
15--30~Å      & 33965  & 0 & 0    & 914.985    & 87.474  \\
более 30~Å    & 193168 & 0 & 2    & 875.171    & 16.174  \\
\bottomrule
\end{tabular}
\end{center}

\smallskip
\noindent Как читать таблицу:
\begin{itemize}[leftmargin=1.4em,itemsep=2pt,topsep=2pt]
\item \textbf{максимум дефектности} — в самой ближней оболочке 0--5~Å: здесь сосредоточены 2894 атома с дефектным окружением и 6 из 7 атомов с гексагональным окружением;
\item \textbf{максимум радиального эквивалентного напряжения по Мизесу} — там же, в оболочке 0--5~Å: 368.194~МПа;
\item уже дальше 15~Å от границы дефектность практически \textbf{исчезает} (в оболочке 15--30~Å — ноль дефектных и ноль гексагональных атомов);
\item дальняя матрица (более 30~Å) остаётся почти идеальной: на 193168 атомов приходится лишь 2 дефектных атома, а эквивалентное напряжение по Мизесу падает до 16.174~МПа.
\end{itemize}
Таким образом, и дефектность, и сдвиговая нагрузка максимальны на самой границе включения и быстро спадают вглубь матрицы.

%======================= 9 =======================
\section{Физическая интерпретация}
\begin{enumerate}[leftmargin=1.6em,itemsep=3pt,topsep=2pt]
\item Магнитострикционное включение Fe$_4$Al$_{13}$ \textbf{способно передавать напряжение} в алюминиевую матрицу. Это видно по радиальному профилю напряжений: у границы возникает заметная сдвиговая нагрузка.
\item Наиболее нагруженная и наиболее дефектная область — \textbf{граница включение--матрица} (оболочка 0--5~Å).
\item Локальное эквивалентное напряжение по Мизесу у интерфейса достигает \textbf{368.194~МПа} (расчёт на 254055 атомов), что \textbf{выше} оценочного предела текучести алюминия около 120~МПа. Это создаёт физические условия для зарождения дефекта.
\item В одном из валидных расчётов (510375 атомов) такой зародыш дефекта действительно обнаружен — короткий зародыш частичной дислокации (1 сегмент, 8.47~Å, тип $\tfrac{1}{6}\langle 112\rangle$), причём только при наличии деформации включения.
\item В длинном расчёте (254055 атомов, 120000 шагов) сформировалась \textbf{локальная дефектная область} у границы (4265 дефектных атомов), но \textbf{без оформленной дислокационной линии}.
\item Итог: гипотеза о локальном механическом воздействии включения на матрицу \textbf{подтверждается}, но \textbf{развитая пластичность пока не доказана} — эффект остаётся локальным и приграничным.
\end{enumerate}

\smallskip
\noindent\textbf{Сравнение двух валидных расчётов.}
\begin{itemize}[leftmargin=1.4em,itemsep=2pt,topsep=2pt]
\item Расчёт на 510375 атомов дал один короткий дислокационный сегмент длиной 8.47~Å.
\item Расчёт на 254055 атомов за 120000 шагов не дал оформленной дислокационной линии (0 сегментов).
\item Оба расчёта \textbf{согласуются в главном}: эффект локальный, возникает у границы включения и не превращается в развитую пластическую зону.
\item Разница (один короткий сегмент против его отсутствия) объясняется размером системы, условиями зарождения дефекта, статистической чувствительностью анализа и тем, что дислокация находится \textbf{на самом пороге появления}: при таком слабом физическом воздействии небольшие изменения объёма, времени и локального теплового шума решают, успеет ли оформиться короткий сегмент.
\end{itemize}

%======================= 10 =======================
\section{Ограничения}
\begin{itemize}[leftmargin=1.4em,itemsep=2pt,topsep=2pt]
\item Локальные напряжения рассчитаны как \textbf{приближённая оценка} по вириальному тензору и оценочному объёму зон; абсолютные значения локальных напряжений нужно трактовать осторожно. Достоверны прежде всего сравнения «зона к зоне» и «физический случай против контрольного», а не точные абсолютные числа в мегапаскалях.
\item Дислокационный анализ \textbf{чувствителен} к размеру расчётной ячейки, времени расчёта, температуре, порогу распознавания дефектов и локальному тепловому шуму.
\item Отсутствие оформленной дислокационной линии в одном расчёте \textbf{не отменяет} факт локальной дефектности у границы.
\item Для строгого доказательства развитой пластичности нужны дополнительные подтверждающие расчёты или прямая визуальная проверка атомной конфигурации.
\end{itemize}

%======================= 11 =======================
\section{Что уже доказано}
\begin{itemize}[leftmargin=1.4em,itemsep=2pt,topsep=2pt]
\item Расчётная часть доведена до \textbf{физически валидного} состояния: стабильная температура около 300~К, отсутствие перегрева и численного разрушения, штатное завершение, 120000 шагов в длинном расчёте.
\item Модель \textbf{соответствует последней постановке} преподавателя: однородная алюминиевая матрица без границы зерна и без вакансий, одно включение Fe$_4$Al$_{13}$, воздействие — собственная деформация включения по оси $Z$ как эквивалент магнитострикции.
\item Включение \textbf{передаёт напряжение} в матрицу; у границы (0--5~Å) локальное эквивалентное напряжение по Мизесу (368.194~МПа) превышает оценочный предел текучести алюминия (около 120~МПа).
\item У границы возникает \textbf{локальная дефектная перестройка} (4265 дефектных атомов в матрице), при этом за пределами ближней интерфейсной оболочки — только 2 дефектных атома, то есть эффект приграничный.
\item В расчёте на 510375 атомов обнаружен \textbf{короткий зародыш частичной дислокации} (1 сегмент, 8.47~Å, $\tfrac{1}{6}\langle 112\rangle$), отсутствующий в контрольном расчёте без деформации.
\end{itemize}

%======================= 12 =======================
\section{Что пока не доказано}
\begin{itemize}[leftmargin=1.4em,itemsep=2pt,topsep=2pt]
\item \textbf{Развитая пластическая зона} в объёме матрицы — не сформирована.
\item \textbf{Развитая дислокационная сеть} — не сформирована.
\item В длинном расчёте на 254055 атомов оформленной дислокационной линии нет (0 сегментов).
\item \textbf{Воспроизводимость} зародыша частичной дислокации не установлена: пока это одно наблюдение (на расчёте 510375 атомов).
\item Абсолютные значения локальных напряжений остаются \textbf{приближёнными} и не могут служить точной количественной мерой.
\end{itemize}

%======================= 13 =======================
\section{Рекомендованный следующий шаг}
\begin{itemize}[leftmargin=1.4em,itemsep=2pt,topsep=2pt]
\item Не плодить большое число расчётных случаев.
\item Научно рациональный следующий шаг — \textbf{один} увеличенный расчёт с тем же физическим параметром деформации 0.001942, но в большем объёме, чтобы проверить, \textbf{повторяется ли} короткий зародыш частичной дислокации у границы.
\item В качестве альтернативы — один контрольный расчёт без деформации включения на такой же геометрии (на системе 510375 атомов такой контроль уже выполнен и дал ноль дислокаций; при увеличении объёма его имеет смысл повторить как опору для сравнения).
\item \textbf{Не делать полный перебор параметров} до обсуждения с преподавателем.
\item Сначала показать настоящий отчёт, получить комментарии и только потом принимать решение о расчёте на больший размер системы (например, около 700 тысяч атомов).
\end{itemize}

\smallskip
\noindent\textit{Примечание. Никакие новые расчёты в рамках подготовки этого отчёта не запускались. Любой следующий запуск должен выполняться только по явному решению пользователя и преподавателя: автоматический контроль текущего расчёта помечен как «требует ручной проверки», а большие расчёты (250k / 500k / 700k) намеренно отключены до отдельного разрешения.}

%======================= 14 =======================
\section{Приложение: исходные числовые таблицы и соответствие терминов}

\subsection{Сводка длинного расчёта на 254055 атомов}
\begin{center}
\small
\setlength{\tabcolsep}{6pt}
\renewcommand{\arraystretch}{1.2}
\begin{tabular}{p{67mm} p{52mm} p{38mm}}
\toprule
\hh{Величина} & \hh{Значение} & \hh{Источник} \\
\midrule
Случай & E3\_phys001942\_250k\_120k & final\_summary \\
Число атомов (факт) & 254055 & final\_summary \\
Расчётная ячейка & 14.175 $\times$ 14.175 $\times$ 21.06 нм & final\_summary \\
Деформация включения & 0.001942 & analysis \\
Шагов выполнено & 120000 / 120000 & final\_summary \\
Финальная температура & 288.39766 К & production\_summary \\
Максимальная температура & 291.552 К & final\_summary \\
Дислокационных сегментов & 0 & analysis \\
Суммарная длина дислокаций & 0.0 Å & analysis \\
Плотность дислокаций & 0.0 м$^{-2}$ & analysis \\
Атомов матрицы & 243669 & analysis \\
Атомов с гранецентр. структурой & 239404 (98.25 \%) & analysis \\
Атомов с гексагон. окружением & 7 & analysis \\
Атомов с дефектным окружением & 4258 & analysis \\
Дефектных атомов в матрице всего & 4265 & final\_summary \\
Дефектных атомов вне оболочки & 2 & analysis \\
\bottomrule
\end{tabular}
\end{center}

\subsection{Сводка расчёта на 510375 атомов}
\begin{center}
\small
\setlength{\tabcolsep}{6pt}
\renewcommand{\arraystretch}{1.2}
\begin{tabular}{p{60mm} r r}
\toprule
\hh{Величина} & \shortstack[r]{\hh{Контрольный}\\\hh{(деформация 0)}} & \shortstack[r]{\hh{Физический}\\\hh{(0.001942)}} \\
\midrule
Число атомов & 510375 & 510375 \\
Максимальная температура, К & 291.5448 & 291.98355 \\
Дислокационных сегментов & 0 & 1 \\
Суммарная длина дислокаций, Å & 0.0 & 8.47 \\
Тип вектора смещения & --- & $\tfrac{1}{6}\langle 112\rangle$ \\
\bottomrule
\end{tabular}
\end{center}
\noindent\footnotesize Источник: \code{stageE\_v2\_analysis\_summary.json}, \code{stageE\_v2\_boundary\_dislocation\_report.md}.\normalsize

\subsection{Соответствие русских терминов и общепринятых названий методов}
\begin{center}
\small
\setlength{\tabcolsep}{6pt}
\renewcommand{\arraystretch}{1.2}
\begin{tabular}{p{82mm} p{72mm}}
\toprule
\hh{Русский термин в отчёте} & \hh{Общепринятое название / сокращение} \\
\midrule
Анализ дислокаций по атомной конфигурации & Dislocation Extraction Algorithm (DXA) \\
Анализ локального кристаллического окружения & Common Neighbor Analysis (CNA) \\
Сопоставление структуры с шаблонами решёток & Polyhedral Template Matching (PTM) \\
Окружение гексагональной плотноупакованной структуры & hexagonal close-packed (HCP) \\
Нераспознанное / дефектное окружение & OTHER \\
Эквивалентное напряжение по Мизесу & von Mises stress \\
Кубическая гранецентрированная структура & face-centered cubic (FCC), ГЦК \\
\bottomrule
\end{tabular}
\end{center}

\subsection{Технические источники данных}
\footnotesize
\noindent Длинный расчёт: \code{runs/stageE\_250k\_single\_physical\_longrun/20260623-205439}
\begin{itemize}[leftmargin=1.4em,itemsep=1pt,topsep=1pt]
\item \code{stageE\_250k\_final\_summary.json}; \code{.../production/analysis.json}
\item \code{stageE\_250k\_dxa\_summary.md}, \code{stageE\_250k\_physics\_verdict.md}, \code{stageE\_250k\_stress\_transfer\_report.md}, \code{stageE\_250k\_longrun\_status.md}
\item \code{tables/defect\_summary.csv}, \code{tables/production\_summary.csv}, \code{summaries/E3\_250k\_longrun\_report.md}
\end{itemize}
\noindent Расчёт на 510375 атомов: \code{runs/stageE\_homogeneous\_inclusion\_scaleup\_v2/20260622-224433}
\begin{itemize}[leftmargin=1.4em,itemsep=1pt,topsep=1pt]
\item \code{stageE\_v2\_analysis\_summary.json}, \code{stageE\_v2\_boundary\_dislocation\_report.md}, \code{stageE\_v2\_stress\_transfer\_report.md}, \code{stageE\_v2\_physics\_verdict.md}, \code{stageE\_v2\_failure\_or\_success\_report.md}
\end{itemize}
\noindent Программный комплекс: LAMMPS (KOKKOS/CUDA), потенциал MEAM (трек \code{MEAM\_Jelinek\_2012}), температура 300~К, шаг 0.0002~пс.
\normalsize

\end{document}
